\section{Step 55 Reconstruct scalar magnitudes and signs}\label{h263-quant-s05}
Read each ordinary level from the private record. For zero, store zero directly.
For a nonzero value, calculate $Q(2|L|+1)$, subtract one if $Q$ is even, then
restore the level's sign. Clip the signed result to $[-2048,2047]$ and store it
in the private reconstruction array.

For INTRA DC, translate code 255 back to selected $d=128$, then multiply $d$ by
eight. This result bypasses the ordinary formula. You may process a placeholder
at index zero in a common loop if you replace it correctly before publication;
the public record must contain the special result.

\paragraph{Observe} At $Q=8$, levels one and minus one reconstruct 23 and
minus 23. Level zero reconstructs zero, not seven.
\Needspace{8\baselineskip}
\paragraph{Check} When is the even-parameter subtraction applied?
\begin{itemize}
\item[A.] After adding a predicted pixel.
\item[B.] To the positive nonzero magnitude before restoring the sign.
\item[C.] To the final pointer.
\end{itemize}
